\documentclass[11pt,a4paper]{article}
\usepackage[margin=1in]{geometry}
\usepackage[T1]{fontenc}
\usepackage[utf8]{inputenc}
\usepackage{lmodern}
\usepackage{microtype}
\usepackage{booktabs}
\usepackage{longtable}
\usepackage{enumitem}
\usepackage{xcolor}
\usepackage{hyperref}
\usepackage{listings}
\usepackage{fancyhdr}
\usepackage{titlesec}
\usepackage{parskip}
\usepackage{amsmath}
\definecolor{pvzgreen}{HTML}{447A36}
\definecolor{pvzdark}{HTML}{263525}
\definecolor{codegray}{HTML}{F4F5F2}
\hypersetup{colorlinks=true,linkcolor=pvzgreen,urlcolor=pvzgreen,citecolor=pvzgreen}
\pagestyle{fancy}
\fancyhf{}
\lhead{PvZ Deep Learning}
\rhead{Technical Development Report}
\cfoot{\thepage}
\titleformat{\section}{\Large\bfseries\color{pvzdark}}{\thesection}{0.7em}{}
\titleformat{\subsection}{\large\bfseries\color{pvzgreen}}{\thesubsection}{0.7em}{}
\lstset{basicstyle=\ttfamily\small,backgroundcolor=\color{codegray},frame=single,breaklines=true,columns=fullflexible}
\setlist[itemize]{topsep=3pt,itemsep=2pt}
\title{\textbf{Plants vs. Zombies Deep Learning}\\\large Technical Development Report}
\author{b3d012}
\date{4 September 2026}
\begin{document}
\maketitle
\begin{abstract}
This report documents the engineering work completed before the machine-learning phase of the \emph{Plants vs. Zombies Deep Learning} project. The project interfaces with the original Windows \emph{Plants vs. Zombies: Game of the Year Edition} process: a read-only memory subsystem reconstructs a structured game state, a placement subsystem converts that state into semantic legality information, and a controller performs ordinary Windows mouse input. Phases 1--3 establish the sensing, actuation, and Environment v1 foundation required by future learning experiments. A subsequent production-runtime layer adds process and window lifecycle ownership, conservative phase/health detection, focus and pause policy, fail-closed action gating, and operator diagnostics without revising the frozen contracts. The report records the research path, reverse-engineering methodology, versioned interfaces, validation methods, known limitations, and the transition to Phase 4.
\end{abstract}
\tableofcontents
\newpage
\section{Project Goal and Design Philosophy}
The long-term goal is to build an AI agent that can play \emph{Plants vs. Zombies} strategically using the real game as its environment. The intended research value is not simply to automate clicks, but to create a reproducible perception--decision--action pipeline suitable for reinforcement learning or another sequential decision-making method.
\begin{itemize}
\item \textbf{Use the original game.} The project targets the Windows GOTY release rather than a reimplemented clone, preserving the actual game dynamics.
\item \textbf{Read internal state directly.} Memory reading gives precise HP, cooldown, wave, type, and coordinate information. Computer vision can still be introduced later as an auxiliary experiment.
\item \textbf{Do not use memory writes for normal agent actions.} Observation is read-only; Controller v1 uses ordinary mouse input.
\item \textbf{Build in phases.} Reader, legality, and controller capabilities were completed independently before a learning loop is introduced.
\end{itemize}
\section{Target Version and Reference Material}
\subsection{Version target}
The reader is pinned to \textbf{PvZ GOTY 1.2.0.1073}. Version-specific offsets are centralized in \texttt{pvz\_reader/versions.py}, preventing build-specific assumptions from spreading through the reader.
\subsection{pvztoolkit as a research reference}
The repository includes \texttt{lmintlcx/pvztoolkit} as a Git submodule under \texttt{references/pvztoolkit}. It was used as a reverse-engineering reference rather than as the AI project's runtime engine. Its source exposed useful pointer chains, structure sizes, seed-bank behavior, and enumerations that could then be checked against the target executable. Seed-bank research, for example, used the toolkit's \texttt{0x50} slot size and its normal/imitater seed fields as hypotheses before live validation.
\subsection{Validation methodology}
Known offsets from other builds were never assumed correct. The practical method became: locate a candidate from source/reference material, inspect the target process, manipulate a single in-game variable when possible, compare snapshots or watch candidate addresses, confirm type/stride/lifetime semantics, encode the result in \texttt{versions.py}, and then expose it through typed state plus focused tests.
\section{Research and Reverse-Engineering Workflow}
The investigation combined C++ source inspection, PowerShell searching, small Python memory-inspection programs, address watches, memory traces, and targeted live experiments. When a preferred command-line search utility was unavailable, PowerShell tools such as \texttt{Get-ChildItem} and \texttt{Select-String} were used to trace relevant fields through the reference source.

Temporary artifacts such as zombie, wave, seed-slot, seed-readiness, pickup, and projectile research traces were intentionally kept while discovery was active. Their durable conclusions are now represented by the version table, implementation, tests, and this report, allowing those raw dumps to be removed from the portfolio-facing working tree.
\subsection{Process and memory access}
The reader attaches to \texttt{PlantsVsZombies.exe} and uses ordinary Win32 process-memory reads. The target build's root chain begins with the LawnApp pointer at \texttt{0x729670}, then the Board offset \texttt{0x868}. The design remains external to the game and avoids DLL injection.
\section{Phase 1: Game-State Reader}
Phase 1 converted the live game process into a coherent machine-readable observation rather than a collection of isolated addresses.
\subsection{Global board state}
The reader exposes sun, game clock, scene, adventure level, and pause state. Key Board offsets for 1.2.0.1073 are sun \texttt{0x5578}, game clock \texttt{0x5580}, scene \texttt{0x5564}, adventure level \texttt{0x5568}, and pause state \texttt{0x017C}.
\subsection{Plants}
The final plant layout uses a \texttt{0x14C} structure stride and captures grid row/column, plant type, state, current/max HP, imitater status, and dead/squished/asleep flags. This makes plant occupancy and condition explicit to downstream logic.
\subsection{Zombies}
The zombie layout uses a \texttt{0x168} stride. The reader captures row, type, status, continuous position, biting state, slow/stun/freeze timers, hypnotized state, body HP, armor HP, maxima, and dead state. This provides tactical information that would be difficult to recover reliably from screenshots alone.
\subsection{Seed bank and cooldown state}
Seed research established the seed-bank pointer, slot count, \texttt{0x50} slot stride, seed type, imitater type, cooldown elapsed/total fields, readiness/cooling flags, and use counter. Together with plant costs this supports derived ready/affordable/actionable state.
\subsection{Wave state}
Wave timing was live validated instead of inferred from UI graphics. The reader exposes total/current/refreshed wave indices, refresh HP, current-wave HP, normal countdown, initial countdown, and huge-wave countdown fields.
\subsection{Lawn mowers}
Mower state includes row, X/Y, internal state, dead/visible status, type, and object ID, allowing the future policy to distinguish lanes that retain an emergency defense from lanes that do not.
\subsection{Pickups}
Pickups required a larger reverse-engineering pass involving board traces, structure traces, version traces, and code dumps. The validated target layout uses Board offset \texttt{0xFC}, capacity \texttt{0x100}, stride \texttt{0xD8}, and fields for X/Y coordinates, collection status, timer, and type. This later enabled Controller v1 to click exact observed pickup positions.
\subsection{Grid items}
Grid-item observation records type, row, column, and dead state using a validated \texttt{0xEC} stride. These objects feed placement legality because some act as board blockers or special supports.
\subsection{Projectiles}
Projectile research established a \texttt{0x94} structure and exposes row, position, type, collision eligibility, and object ID. Projectile observations can be omitted from an initial baseline while remaining available for richer tactical models and ablations.
\subsection{GameState v1 freeze}
The Phase 1 result is \texttt{GameState} schema version 1 containing global state, plant/zombie capacities, wave data, plants, zombies, seeds, mowers, pickups, projectiles, and grid items. Placement legality and controller/action state are intentionally not embedded inside GameState; it remains an observation of the game.
\section{Placement Legality and Action Masks}
\texttt{pvz\_reader/placement.py} adds a deterministic legality layer before the agent is allowed to click the board. It handles ordinary grass as well as roof Flower Pots, Lily Pads, water plants, Grave Buster, Pumpkin, Coffee Bean, and upgrade plants including Gatling Pea, Twin Sunflower, Gloom-shroom, Cattail, Winter Melon, Gold Magnet, Spikerock, and Cob Cannon.

Checks cover bounds, seed existence/actionability, cooldown, affordability, scene/terrain, graves/craters, tile occupancy, support-plant behavior, and upgrade prerequisites. The layer returns machine-readable reason codes and can build a placement mask indexed by seed slot, row, and column. This is directly reusable as an invalid-action mask in Phase 3.
\section{Phase 2: Controller v1}
Controller v1 completes the action half of the loop. Higher layers request semantic actions; the controller translates those requests into safe physical input.
\subsection{Logical coordinate system}
PvZ uses an 800 by 600 logical client. \texttt{pvz\_controller/coordinates.py} defines lawn tile, seed packet, shovel, and pickup click points in logical coordinates and scales them to the actual client size. Standard board geometry uses nine columns and up to six rows. Shovel placement accounts for the seed bank widening as more packet slots are unlocked.
The compatible v0.2.1 patch corrected two physical-input assumptions without
changing Controller v1 semantics. On the supported Adventure 1-7 client, the
six-packet seed bank was calibrated as $(120 + 59s, 43)$ for slot $s$; the
count-dependent spacing rule remains derived from reference source. Normal
five-lane lawn input was calibrated as $(80 + 80c, 130 + 100r)$ for column
$c$ and row $r$. In particular, $x=120$ is a lawn-cell boundary interpretation
rather than the validated robust first-column input point. All six packet
selections and exact memory-confirmed placements across all five normal lanes
were validated live; pool/fog and roof mappings remain separately tested.
\subsection{Safe Windows input backend}
\texttt{pvz\_controller/windows\_input.py} enumerates visible windows belonging to \texttt{PlantsVsZombies.exe}, prefers the expected title, rejects minimized/zero-sized clients, obtains client bounds, converts logical client positions into screen positions, focuses the game, moves the cursor, and emits one ordinary left click through \texttt{SendInput}. An injectable Win32 API seam allows this logic to be tested without clicking the real desktop.
\subsection{Semantic actions}
\texttt{PvZController} exposes three primary Controller v1 actions: collecting an observed pickup, planting a selected seed at a board tile, and shoveling a board tile. Each action returns an \texttt{ActionResult}. Planting validates game state, seed readiness/affordability, deterministic placement legality, and both click targets before starting the two-click sequence. Observation-based helpers verify whether a pickup disappeared, a plant appeared, or a plant was removed.
\subsection{Safety properties}
The frozen controller rejects unusable/paused states, bounds-checks logical and scaled points, validates target-window availability and focus, prevalidates multi-click planting before selecting a packet, makes OS input failures explicit, and supports postcondition verification from the reader.
\section{Testing and Validation}
The project separates deterministic tests from explicitly named live tools so normal validation cannot accidentally manipulate a running game.
\subsection{Offline tests}
\begin{longtable}{p{0.34\textwidth}p{0.58\textwidth}}
\toprule Test & Responsibility \\
\midrule\endhead
\texttt{test\_game\_state\_schema.py} & Frozen GameState contract. \\
\texttt{test\_pickups.py} & Pickup decoding/data behavior. \\
\texttt{test\_projectiles.py} & Projectile decoding/data behavior. \\
\texttt{test\_placement.py} & Ordinary and special placement legality, reasons, and masks. \\
\texttt{test\_controller\_schema.py} & Public Controller v1 API contract. \\
\texttt{test\_controller\_coordinates.py} & Tile/seed/shovel/pickup coordinate transforms, scaling, bounds. \\
\texttt{test\_controller\_windows\_input.py} & Window targeting, focus, scaling, and failure behavior through the injectable Win32 seam. \\
\texttt{test\_controller.py} & Semantic action preconditions, click sequences, failures, and post-observation helpers. \\
\bottomrule
\end{longtable}
\subsection{Live validation}
Live validation tools include \texttt{live\_test\_placement.py}, \texttt{live\_test\_placement\_mask.py}, \texttt{live\_test\_controller\_pickup.py}, \texttt{live\_test\_controller\_plant.py}, and \texttt{live\_test\_controller\_shovel.py}. A useful lesson occurred during planting validation: an apparent failure was caused by testing a level where only the three middle lawn rows were active. Repeating the test on a valid row confirmed the planting path. This reinforced the need to validate both internal state and scene geometry.
\subsection{Phase freeze}
Controller v1 was frozen in commit \texttt{7777914} (\texttt{Freeze Controller v1 API}), and Phase 2 was subsequently marked complete by commit \texttt{1889b728} (\texttt{Phase 2 Finished}). The freeze gives future model code a stable observation/action contract.
\section{Repository Cleanup Rationale}
Durable assets include \texttt{pvz\_reader}, \texttt{pvz\_controller}, placement logic, deterministic tests, live-validation scripts, reusable inspection tools, and the \texttt{references/pvztoolkit} submodule. Generated \texttt{\_\_pycache\_\_} data, raw root-level traces/code dumps, and the proprietary local game installation do not belong in the portfolio-facing working tree once their findings are documented.

Deleting the game directory from the current tree is not the same as deleting its historical blobs. Because the installation was previously committed, a separate \texttt{git filter-repo} or equivalent history rewrite is recommended before broadly promoting the repository if the objective is to remove redistributed proprietary binaries from all reachable history.
\section{Status at the End of Phase 2}
The low-level loop is now:
\[
\text{PvZ process}\rightarrow\text{GameState v1}\rightarrow\text{legality/masks}\rightarrow\text{semantic action}\rightarrow\text{Windows input}\rightarrow\text{PvZ process}.
\]
The next work should therefore focus on the learning-environment contract rather than more low-level automation.
\section{Recommended Phase 3: Learning-Environment Bridge}
\subsection{Objective}
Phase 3 should first make the real game look like a standard sequential decision environment, preferably through a Gymnasium-style interface, rather than immediately training a large neural network.
\subsection{Observation encoder}
GameState contains variable-length entity lists and Python objects. Add a deterministic encoder producing fixed, normalized numeric features: global sun/time/wave data; per-tile plant type/HP/state/support channels; spatial or lane-binned zombie type/position/HP/status; seed type/cost/readiness/cooldown; mower state; and optional projectile/pickup channels. Log the raw state alongside encoded observations so trajectories survive future encoder changes.
\subsection{Action space and masking}
Use semantic discrete actions such as \((\text{seed slot},\text{row},\text{column})\), plus no-op/wait and optionally shovel actions. Reuse the existing placement mask to eliminate impossible choices. Pickup collection is primarily mechanical, so an initial agent can use an automatic pickup helper between strategic decisions rather than forcing the policy to learn cursor housekeeping.
\subsection{Step contract}
One environment step should: observe, encode, build the action mask, accept an action, execute through Controller v1, advance for a documented interval/event, observe again, verify action outcome, compute reward and terminal flags, and log the transition.
\subsection{Episode reset}
The first environment can target one reproducible level and permit a manually prepared start while the step loop is validated. Reliable automated menu/reset navigation can be added later instead of destabilizing Controller v1 now.
\subsection{Transition logging}
Persist timestamp, episode ID, raw state, encoded observation, mask, requested action, action result, next state, reward, terminated, and truncated. This produces debuggable trajectories and enables later offline/imitation-learning experiments.
\subsection{Reward design}
Keep terminal win/loss dominant. Add only small, auditable shaping terms based on observed transitions (for example wave progress or immediate threat change) and test them for reward hacking.
\subsection{Definition of done}
Freeze Phase 3 when encoding and masks are deterministic, invalid actions are masked/rejected, action results reconcile with the next observation, step timing and terminal rules are documented, trajectories round-trip to disk, and a random/scripted baseline can run repeated steps or episodes without corrupting the interface. Only then move to the main learned policy, with PPO plus invalid-action masking a sensible first baseline.
\section{Phase 3: Environment v1}
\subsection{Objective and boundary}
Phase 3 converts the frozen reader and controller into a deterministic, library-independent reinforcement-learning environment bridge. It does not implement a learned policy, Gymnasium adapter, menu automation, or a game clone. The environment continues to observe the real target process read-only and requests normal Controller v1 mouse input only through semantic actions.

\subsection{Observation v1}
Observation v1 is a fixed \texttt{float32} vector of shape \texttt{(5534,)}. It preserves a separate raw \texttt{GameState v1} alongside encoded data. The representation includes normalized global and wave features, Adventure level, a six-by-nine plant map, ten seed-bank slots, mower state, and up to five closest-to-house zombie slots per lane. Bounded per-lane live and overflow aggregates retain information when more than five zombies are present. Pickups, projectiles, and grid items are deliberately deferred from the encoded vector; raw grid items still feed deterministic placement legality. Feature order, normalization, and schema version 1 are frozen.

\subsection{Action v1 and deterministic legality}
Action v1 contains 541 indices: \texttt{WAIT} at index zero and 540 \texttt{PLANT(seed\_slot,row,column)} actions in seed-slot, row, then column order. The action mask delegates every placement decision to the existing placement layer rather than asking a policy to rediscover terrain, occupancy, prerequisites, affordability, or cooldown rules. Frozen \texttt{GameState v1} does not authoritatively expose early-level inactive rows; each episode therefore supplies an explicit immutable six-boolean active-row configuration. Shovel and pickup actions remain deferred.

\subsection{Environment v1 step and reconciliation contract}
\texttt{PvZEnvironment} is constructed from reader, controller, encoder, timing, and persistence seams. An explicit \texttt{reset(EpisodeConfig)} adopts a manually prepared, available, unpaused game state and returns a raw snapshot, encoded observation, and action mask. The lifecycle is \texttt{UNINITIALIZED}, \texttt{ACTIVE}, \texttt{TERMINATED}, or \texttt{TRUNCATED}; no step is permitted before reset or after an outcome. A legal WAIT or PLANT advances once for the configured strategic interval and reads a post-step state. If an issued PLANT is not yet visible, the environment may take additional bounded, read-only reconciliation reads (0.75 seconds maximum at 0.05-second intervals by default); this verification window is explicitly separate from agent step advancement and never reissues input. WAIT still takes exactly one post-step read. \texttt{StepTiming} records reconciliation polls separately. Plant outcomes distinguish observed placement, controller failure, missing postcondition, and unavailable post-state. These seams make offline tests deterministic without desktop input.

\subsection{Transition records and reward outcomes}
Transition schema version 2 writes append-only UTF-8 JSONL records containing raw before/after state dictionaries, encoded arrays with dtype and shape, masks, action, controller result, reconciliation, timing, reward, terminal/truncation flags and reason, reward components, and reward schema/spec identity. Episode IDs and step indexes are reset-owned; rejected attempts are recorded and reset itself does not create a synthetic transition.

Reward schema version 1 uses an inspectable frozen \texttt{RewardSpec}. A detected win receives +1.0 and a detected loss -1.0. The only positive shaping term is +0.01 per positive spawned-wave delta. Rejected actions and technical failures receive small diagnostic penalties; WAIT, successful planting, sun spending, and unchanged progress receive no activity reward. Terminal reward suppresses shaping. This delta-based design reduces obvious reward farming. Natural terminal outcomes remain separate from externally imposed truncations such as a maximum step horizon or repeated unavailable state. No authoritative win/loss signal exists in GameState v1, so a validated injected terminal detector is required; the default detector deliberately makes no guess.

\subsection{Episode lifecycle and baselines}
\texttt{EpisodeConfig} freezes episode identity, active rows, timing and truncation limits, Reward v1 configuration, optional detector, and caller-provided metadata. Reset is deliberately manual: the operator prepares the selected level; the environment never navigates menus, seed selection, dialogs, retries, or Adventure progression. A stateful detector may receive a reset callback. The same transition sink can span episodes while retaining distinct IDs and indexes starting at zero.

The evaluation API provides a seeded random valid-action policy, a small scripted economy/threat policy, decision reasons, \texttt{run\_episode}, \texttt{EpisodeResult}, and deterministic multi-episode summaries. The random baseline uses the encoded observation and mask. The scripted baseline may inspect the current structured snapshot to select a legal Sunflower economy action or Peashooter in a threatened lane; it is an auditable engineering baseline rather than a neural-policy architecture. Both execute only Action v1 indices through Environment v1.

\subsection{Frozen interface summary and validation}
Environment v1 freezes Observation schema 1 and shape \texttt{(5534,)}, Action schema 1 and 541-action indexing, Environment schema 1, Reward schema 1 default weights, and Transition schema 2. The \texttt{environment\_contract()} helper returns this deterministic metadata for future checkpoints and evaluation artifacts. The Phase 3 offline suite runs on Windows CI and covers deterministic observation encoding, masks, controller seams, environment lifecycle, rewards, JSONL round trips, baselines, and frozen contract metadata.

An explicit live tool, \texttt{tools/live\_run\_environment.py}, supports dry-run observation/mask/policy validation by default and requires \texttt{--execute} before it can issue normal mouse input. It requires explicit active rows and a bounded horizon. Dry-run validation confirmed reset, the \texttt{(5534,)} observation, the \texttt{(541,)} action mask, legal semantic action selection, and zero Controller input. JSONL round-trip validation preserved observation and mask shapes together with episode, step, reward, and truncation data.

End-to-end validation against the real client then ran both the scripted heuristic and a seeded random-valid-action baseline for ten strategic steps on a normal five-row lawn. Each run executed eight WAIT actions and two PLANT actions, reached the configured \texttt{MAX\_STEPS} truncation, and sent no rejected or masked action to the Controller. In each run, an initial transient Windows focus failure was correctly surfaced as \texttt{CONTROLLER\_FAILED}; the following legal Sunflower action reconciled as \texttt{PLANT\_OBSERVED}. The remaining WAIT actions reconciled as \texttt{WAIT\_ADVANCED}. The heuristic run recorded cumulative reward $-0.005$; the seeded random run recorded $+0.005$, including one $+0.01$ Reward v1 wave-progress transition. Transition logging was validated throughout. Inactive-row blocking remains a deterministic offline regression test rather than a deliberately forced live click.

The initial live investigation exposed a reconciliation race: the game accepted a legal plant (as indicated by sun expenditure) before the first post-step snapshot made it visible, yielding false \texttt{PLANT\_NOT\_OBSERVED} results. Environment v1 was corrected with bounded, read-only postcondition polling after the one configured strategic advance. Both post-fix live runs recorded zero \texttt{PLANT\_NOT\_OBSERVED} outcomes. The observed focus failure is an explicit desktop/controller runtime condition, not a policy or environment correctness failure. Generated logs and trajectories are ignored.

\subsection{Known limitations and transition to Phase 4}
Environment v1 deliberately has no automated reset or pickup collection, no authoritative natural terminal detector, no Gymnasium adapter, and no learned model. Manual level preparation and explicit active rows are required for live episodes. Natural win/loss remains unvalidated because GameState v1 has no authoritative detector; the live runs therefore use configured max-step truncation. These limitations are documented constraints rather than hidden heuristics. Phase 4 may introduce deep reinforcement-learning experiments only while preserving the frozen Phase 1--3 public contracts or explicitly versioning any justified change.

\section{Production Runtime Infrastructure}
\subsection{Motivation and requirements}
Environment v1 provides a deterministic strategic step contract, but long-running training and operator tooling also require ownership of operating-system state that does not belong in GameState: process lifetime, PID changes, window identity, foreground policy, attachment health, and safe pause control. The runtime milestone therefore establishes a production boundary beneath Environment v1. Its priorities are fail-closed input, backward compatibility, explicit diagnostics, controlled recovery, and dependency-light operation. It does not add rewards, menu automation, a Gymnasium API, or learned behavior.

\subsection{Architecture and compatibility boundary}
The new \texttt{pvz\_runtime} package contains three layers. \texttt{PvZSession} owns the PID-bound process, memory reader, game-state reader, and matching window. \texttt{GamePhaseDetector} and \texttt{EnvironmentHealth} interpret conservative state and safety evidence. \texttt{PvZRuntime} serializes observation and control, applies the watchdog, and delegates accepted semantic actions to Controller v1. Reader and plant-controller adapters place this infrastructure beneath the existing \texttt{pvz\_env.PvZEnvironment}. This composition preserves GameState v1, Controller v1, Observation v1, Action v1, Environment v1, Reward v1, and Transition schema v2.

\begin{center}
\begin{tabular}{ll}
\toprule
Component & Responsibility \\
\midrule
\texttt{PvZSession} & process discovery, PID attachment, window binding, detach/reattach \\
\texttt{GamePhaseDetector} & conservative semantic phase from validated evidence \\
\texttt{EnvironmentHealth} & separate observation and action readiness with reasons \\
\texttt{PvZRuntime} & focus, pause, freshness, watchdog, action dispatch, snapshots \\
Runtime monitor & operator frontend using the same runtime methods \\
Environment adapters & backward-compatible integration beneath Environment v1 \\
\bottomrule
\end{tabular}
\end{center}

\subsection{Session management and process discovery}
Process discovery accepts only the supported \texttt{PlantsVsZombies.exe} and \texttt{popcapgame1.exe} names and, where available, validates the executable-path basename. Memory attachment is performed by PID rather than by a second independent name lookup. The session owns and closes the memory handle. Each \texttt{ensure\_attached()} call performs at most one recovery attempt: a dead PID invalidates the old reader, releases its handle, discovers the replacement process, and increments a connection generation. There is no unbounded internal reconnect loop.

\subsection{Window discovery and version confidence}
The Windows backend now binds enumeration to both supported executable name and the session's exact PID. This prevents controller targeting from drifting to another same-named process during restart races. Window validity includes existence, non-minimized state, usable client dimensions, and PID agreement. The runtime reports the expected GOTY layout \texttt{1.2.0.1073}; because the repository has no authoritative executable fingerprint, \texttt{version\_verified} deliberately remains false. Coherent Board reads provide useful operational evidence but are not presented as binary-version proof.

\subsection{Focus policy and action gating}
\texttt{FocusMode.MANUAL} requires the PID-bound game window already to be foreground. The runtime refuses an unfocused request, and the input backend independently refuses if focus is lost in the final race before input. \texttt{FocusMode.AUTO} may restore a minimized target and attempt activation through documented Win32 window and input-queue operations: relevant GUI threads are attached temporarily, the target is raised and activated, attachments are removed in a \texttt{finally} path, and the exact foreground HWND is polled for a short bounded interval. Failure still prevents input. Every semantic action is serialized and passes process, reader, Board, state-age, phase, pause, window, and focus gates. Focus is followed by a fresh observation and a second state-gate pass before Controller v1 dispatch. Expected failures return a typed \texttt{RuntimeActionStatus}; Controller v1 results are retained when dispatch occurs.

\subsection{Pause and resume semantics}
\texttt{pause()}, \texttt{resume()}, and \texttt{set\_paused(bool)} are idempotent. The current GameState pause flag is read first; a matching state produces no input. A required transition satisfies the configured focus policy, sends exactly one Escape press to the verified foreground PvZ window, and polls read-only state within explicit timeout and interval bounds. For compatibility with the original client, Escape is encoded as one \texttt{MapVirtualKeyW}-derived scan-code key-down/key-up pair using \texttt{KEYEVENTF\_SCANCODE}. No virtual-key fallback is sent, since two accepted logical presses could immediately reverse the transition. Results distinguish an already satisfied request, a verified change, missing state, focus failure, input failure, and transition timeout; \texttt{CHANGED} is reported only after memory verification.

\subsection{Game-phase detection}
Phase classification uses only process liveness, reader validity, Board presence, \texttt{GameState.paused}, and \texttt{game\_clock}. A coherent Board at clock zero is \texttt{READY}; a later unpaused Board is \texttt{PLAYING}; and the validated pause flag produces \texttt{PAUSED}. With no Board, current evidence cannot reliably distinguish menu, loading, seed selection, result screens, or short transitions, so the honest phase is \texttt{MENU\_OR\_TRANSITION}. Invalid/incomplete evidence is \texttt{UNKNOWN}. Missing-Board display transitions are minimally debounced, while Board health fails on the first missing read. Natural \texttt{LEVEL\_WON} and \texttt{LEVEL\_LOST} values require an independently validated injected provider because frozen GameState v1 has no authoritative terminal field.

\subsection{Watchdog and stale-state protection}
\texttt{EnvironmentHealth} reports process, window, focus, reader attachment, reader validity, controller readiness, Board validity, phase, state age, focus mode, and diagnostic reasons. \texttt{can\_observe} requires a current live attachment, valid reader, Board, and non-stale state. \texttt{can\_act} additionally requires the \texttt{PLAYING} phase, a usable PID-bound window, controller readiness, focus satisfiable under policy, and non-observer mode. \texttt{READY} remains observable but is deliberately not actionable because a clock-zero Board can still be transitional. A fresh read is mandatory for each action, and an observer-only configuration makes all input unavailable while retaining diagnostics.

\subsection{Snapshots, logging, and monitor}
Runtime events use Python's standard \texttt{logging} package at informative levels without logging every successful read. \texttt{RuntimeSnapshot.to\_dict()} produces compact JSON-compatible process/window identity, expected and foreground HWND values, health, phase, level, scene, clock, pause, sun, wave, and entity-count data plus the most recent action, focus, pause, input, and error outcomes. It excludes raw entity arrays and unrelated machine information.

The dependency-light Tk monitor displays connection, PID, window title, expected/foreground HWNDs, focus and focus mode, reader/controller/Board readiness, phase, Adventure level, wave, pause, sun, entity counts, state age, and recent diagnostics. Focus, Pause, Resume, Refresh/Reattach, Snapshot, and Detach controls all invoke \texttt{PvZRuntime}; the UI contains no duplicate controller logic. One worker thread prevents memory/process polling from blocking repainting, while a bounded FIFO preserves accepted operator commands exactly once. Automatic refreshes are coalesced and submitted only while no command is pending or active, so they cannot flood the worker or discard a button request. Each completed operation exposes its status and detail in the UI; shutdown places runtime cleanup after active work.

\subsection{Testing methodology and live validation}
Offline tests inject process discovery, readers, clocks, sleepers, windows, and controllers. Coverage includes process absence, attachment, missing windows, process restart and PID replacement, stale-reader invalidation, detach, conservative phases and debounce, MANUAL/AUTO focus outcomes, temporary Win32 input-queue attachment and cleanup, scan-code input structures, Board/window/pause/observer watchdog refusals, stale-state reporting, idempotent pause/resume, input failure and timeout, typed controller rejection, snapshot diagnostics, Environment v1 adapters, PID-bound window input, monitor result presentation, command preservation during refresh, refresh coalescing, and shutdown without opening a GUI. The complete pre-existing reader, Controller v1, and Environment v1 suites remain part of the same Windows CI run.

The existing Environment v1 gameplay path has already been validated against the real client. Dedicated runtime validation is intentionally separate: \texttt{tools/live\_test\_runtime.py} defaults to read-only diagnostics, requires explicit flags and interactive confirmation for focus or Escape input, and never plants, shovels, collects, or navigates menus. A read-only live pass automatically discovered the running process, bound the matching titled 800 by 600 window to the same PID, read a coherent Board, classified its paused state as \texttt{PAUSED}, and reported observation healthy while action remained unavailable because the client was paused and unfocused. A subsequent AUTO-focus attempt was denied by Windows and correctly failed closed without Escape input.

The first real monitor pass then confirmed correct GUI launch, process/window identity, reader/controller/Board health, live state updates, manual \texttt{PLAYING}/\texttt{PAUSED} tracking, and visible Focus Game behavior. It also exposed a critical integration defect: a single Future represented both periodic refresh and controls, so a click arriving during refresh was silently discarded. Mock tests had correctly established idempotence and fail-closed orchestration, but could not prove monitor command scheduling, Windows foreground acquisition, or the game's acceptance of virtual-key Escape input. The implementation was corrected with the command FIFO, stronger bounded foreground acquisition, scan-code Escape injection, and explicit operation diagnostics described above.

The corrected path was operator-validated against the real client on 4 September 2026. Read-only attachment identified the target process, its PID-bound game window, and a coherent Board while the game was paused. AUTO focus acquired the expected HWND before input; when that identity could not be proved, the runtime returned \texttt{FOCUS\_REQUIRED} and sent no input. A scan-code Escape pause/resume cycle was observed in both memory and the client UI, without duplicate toggles. Repeated pause and resume requests were idempotent. The monitor controls completed and displayed their results, including a deliberate manual-focus failure that left the game unchanged. A restart created a fresh PID/HWND identity without stale-handle reuse. These results freeze the managed runtime as Phase 3.5 and make clear that foreground denial is a Controller/runtime failure mode, not a policy or environment correctness failure.

\subsection{Design tradeoffs, limitations, and Phase 4 readiness}
The implementation favors a small synchronous core and one lock rather than an autonomous watchdog thread or asynchronous framework. Reconnection is demand-driven, avoiding runaway polling and orphan workers. Standard Tk and logging avoid new dependencies. The cost is that process recovery occurs on the next explicit operation rather than instantaneously.

Known limitations are explicit: exact executable version is not fingerprinted; Board absence does not identify a specific application screen; natural win/loss still lacks validated evidence; Windows may deny foreground activation; menu navigation, seed selection, dialogue handling, pickup automation, and level restart remain outside scope. Phase 4 can consume the frozen Environment v1 through runtime adapters without learning code knowing PIDs, handles, HWNDs, focus mechanics, pause toggles, memory addresses, or raw coordinates.

\section{Phase 4 Training Lifecycle Extension Candidate}
The unreleased lifecycle extension retains every frozen v1 interface and adds
orthogonal runtime services. Natural terminal evidence is read from the
supported GOTY application's scene and Board result together with the live
Board's level-complete flag. The provider exposes \texttt{RUNNING},
\texttt{WON}, \texttt{LOST}, or \texttt{UNKNOWN} plus raw evidence. A live
playing Board overrides a retained prior result; ambiguous transitions remain
unknown.

Reset is split into control and proof. A pluggable driver requests restart;
the verifier returns \texttt{RESET\_OK} only for a replaced Board, the same
Adventure level, optional exact seed types, an unpaused near-initial clock,
observable health, and optionally empty entity sets. The default driver
refuses, preventing an unvalidated coordinate sequence from masquerading as
automatic reset. Managed pickups use Controller v1 under the runtime's single
input lock, deduplicate pending identities, and confirm disappearance. Offline
tests cover failure modes; real-client checks remain mandatory before v0.2.0.
An observer-only partial live observation on 4 September 2026 read a coherent
paused Adventure level 7 Board with application scene 3, Board result 0, and a
false completion flag, producing \texttt{RUNNING}. No input was sent. This does
not validate win, loss, reset, or pickup behavior.

\section{Phase 4 Training Lifecycle Release}
Version 0.2.0 retains every frozen v1 interface and adds a narrow lifecycle
service above them. The supported research condition is Adventure 1-7 on the
Windows GOTY 1.2.0.1073 client. Earlier forced-level experiments were unstable
on the target installation, so the condition is prepared normally and checked
through memory.

Natural terminal evidence is separated from GameState v1. A live loss is
identified by the zombies-won scene. A post-Board win retains the application
result. The important live-Board win finding is
\texttt{Board::mLevelAwardSpawned} at offset \texttt{0x5624}: during a visible,
uncollected reward-card state at level 7, wave 20/20 and zero zombies, the
field was true while \texttt{mLevelComplete} was false. Thus a retained
\texttt{mBoardResult} alone is deliberately never accepted as a live-Board win.

The reset driver is intentionally version-pinned to the 800 by 600 client.
Menu, Restart Level, and native Try Again use logical coordinates
\texttt{(739,13)}, \texttt{(400,358)}, and \texttt{(384,369)} respectively,
with a 100 ms cursor-settle delay. The outer verifier requires a distinct
Board, same Adventure level, fresh unpaused clock, expected seeds when
configured, and clean entities. Live validation passed active restart,
known-menu restart, four clean loss resets, three clean reward-pending win
resets without collecting the reward card or entering level 8, and 11/11
managed sun-pickup attempts. Pickup and strategic input remain serialized by
the runtime lock.

\section{Historical Pre-Phase-3 Portfolio Work}
Before training begins: remove/ignore the local game, generated caches and raw traces; consider purging the game from Git history; add a concise root README with architecture/status/setup/report link and a statement that game files are not included; add reproducible Python environment metadata; add Windows CI for non-live tests; and tag the cleaned Phase 2 milestone (for example \texttt{phase-2-controller-v1}).
\section{Conclusion}
Through Phase 3.5, the project has solved the integration problems that normally disappear behind a simulator: version-specific reverse engineering, precise structured observation, rule-aware invalid-action masking, safe actuation of a real desktop application, typed reward and lifecycle semantics, portable trajectory logging, baseline evaluation against the real client, process/window lifecycle management, and fail-closed runtime control. Environment v1 and its managed runtime harness are frozen for the v0.1.0 release. The resulting portfolio narrative is therefore \textbf{reverse engineering $\rightarrow$ structured state estimation $\rightarrow$ rules/action masking $\rightarrow$ safe control $\rightarrow$ Environment v1 $\rightarrow$ managed runtime $\rightarrow$ learned strategy}. Phase 4 can now concentrate on learning experiments without exposing model code to process IDs, Windows handles, focus mechanics, or memory-reader lifecycle.
\end{document}
